\chapter{神经网络的泛化}
\label{chap:neural-network-generalization}

当网络已经能把训练误差降得很低，模型选择还没有结束。两个网络可能在训练集上同样准确，却对未见过的背景、主体或输入变化给出不同预测。增加训练时间、换一种初始化，或者修改数据增强，都可能改变这种差异。我们需要回答的不只是“网络拟合了多少样本”，还包括“它从样本中保留了什么规律，以及这些规律何时能够用于新数据”。

第一卷的学习理论已经区分经验风险与总体风险，并提供复杂度、稳定性和归纳偏置的分析工具。本章将这些认识放回神经网络的实际学习过程：数据增强改变网络接触的经验，模型结构与显式损失规定可选函数和评价偏好，更新、随机扰动与轨迹选择影响实际学得的解，集成则组合多个已学得的预测。这些方法作用位置不同，适用条件也不同。

记参数为$\bm\theta$，训练集为$S=\{(\bm x_i,y_i)\}_{i=1}^n$，未来样本分布为$P$。经验风险与总体风险分别为
\begin{equation}
  \hat R_S(\bm\theta)=\frac1n\sum_{i=1}^n
    \ell(f_{\bm\theta}(\bm x_i),y_i),\qquad
  R_P(\bm\theta)=\mathbb E_{(\bm x,y)\sim P}
    [\ell(f_{\bm\theta}(\bm x),y)]\text{。}
  \label{eq:generalization-risks}
\end{equation}
以下方法通常希望降低未来风险，而不只是缩小$R_P-\hat R_S$：一个同时在训练和未来数据上很差的常量模型，也可能具有很小的风险差。调节增强强度、正则化系数和停止时刻使用独立验证数据，最终测试数据用于评估已确定的方案。

\section{主要挑战}
\label{sec:generalization-challenges}

第\ref{sec:deep-theory-comparison}节分别从参数大小、初始梯度、更新规则和参数分布理解可学习性。这里把实际泛化困难对应到数据是否有代表性、允许哪些函数及算法最终选出哪个解，作为后面方法的选择依据。

\subsection{数据不足与数据质量}
\label{sec:generalization-data-quality}

经验风险只约束已经出现的样本，优化器继续降低它也不能补足训练分布缺失的情形。因此，数据覆盖与质量先决定训练目标是否代表未来任务，再决定增强或正则化应处理什么缺口。若识别器只见过某类物体与固定背景同时出现，就可能将背景当作类别线索；这种规则能够降低训练误差，却未必适用于背景变化后的物体。覆盖不足不是单纯的样本总数问题，还涉及类别、环境、主体和时间范围。

标签噪声又使真实规律与标注误差混在一起。独立、近似随机的标注错误与只在某类样本上发生的系统性错误具有不同影响。增加来自同一主体的重复记录，可能扩大名义样本数，却没有提供同等数量的独立信息。第\ref{chap:statistical-learning-framework}章中的独立同分布假设应按真实采样单位检查，不能把相关窗口或复制样本直接当成独立观察。

\term{同分布泛化}（in-distribution generalization）考察来自与训练采样相同分布的新样本；\term{分布变化}（distribution shift）则涉及未来输入、标签机制或二者关系发生改变。前者的风险控制不能自动保证后者的表现。数据增强可以模拟某些已知变化，但不能凭空覆盖未知机制。对于主体相关、时间相关或场景相关数据，应按预期使用条件划分训练、验证和测试，避免共享信息造成过于乐观的评估。

\subsection{模型容量与归纳偏置}
\label{sec:generalization-capacity}

第\ref{sec:network-learning-problem}节中的模型族与样本量分析说明，训练拟合能力要连同函数范围理解。共享和范数约束通过改变或偏好这个范围减少自由度，其收益取决于是否保留任务需要的规律。大网络能够表示复杂规律，也可能拟合训练数据中缺乏可迁移意义的细节。\term{模型容量}（model capacity）描述模型族能够实现多丰富的函数或判别行为；参数数量只是其中一个尺度。权重范数、层间共享、激活、输入分布和学习算法会改变实际可用的函数范围，不能仅由“参数多于样本”判断最终泛化。

\term{归纳偏置}（inductive bias）是模型与学习方法倾向于采用哪些规律的先验偏好。卷积倾向复用局部模式，循环共享倾向使用同一状态更新，注意力的可见性和位置表示规定哪些交互可以发生。这些偏好与任务吻合时，可以减少从数据中重新学习结构的负担；与任务不符时，也可能排除需要的规律。

第\ref{sec:network-norm-bounds}节说明了网络范数如何进入某些泛化界。界依赖输入、损失和参数约束，常是充分条件或上界；它不意味着只要把某个范数压低，预测风险就一定降低。容量不足导致拟合偏差，约束过弱则可能保留无关自由度，实际选择需要在代表性验证数据上比较。

\subsection{解选择与训练敏感性}
\label{sec:generalization-solution-selection}

第\ref{sec:implicit-bias}节说明，相同目标和函数范围仍可因更新规则不同而选出不同解。泛化控制因此还需要关注训练轨迹，而不只修改网络大小或惩罚系数。初始化、样本抽样、训练顺序和优化器决定学习过程到达哪一个参数位置，因而也可能改变未见数据上的预测。\term{训练敏感性}（training sensitivity）在这里指训练数据或训练随机性的小变化导致所得预测明显变化，不等同于模型对输入扰动的敏感性。

正则化可以排除某些解，提前终止限制进一步拟合，参数平均综合同一轨迹的多个状态，模型集成则组合多个预测。这些方法不必把单模型训练误差继续压低。第\ref{sec:algorithmic-stability}节中的算法稳定性关注更换一个训练样本怎样影响输出损失，它提供了把训练过程与泛化联系起来的途径；某些网络上的经验稳定，并不能替代相应假设下的形式化保证。

\section{数据增强}
\label{sec:generalization-augmentation}

当训练数据没有充分展示任务允许的变化时，数据增强把这些变化写入新的经验目标，要求一个参数模型同时解释原样本及相应变换。它通过任务先验约束预测行为，与第\ref{sec:generalization-capacity}节的归纳偏置相联系；增加训练视图并不增加同等数量的独立样本。

\subsection{基于变换的数据增强}
\label{sec:generalization-transform-augmentation}

变换增强针对已有样本附近的覆盖缺口，将“不应改变答案”或“答案应怎样同步变化”变成损失约束。它改善的是目标覆盖，而非第\ref{sec:training-curvature}节中损失曲率的必然条件。若任务认为某些输入变化不应改变答案，可以在训练中主动展示这些变化。\term{数据增强}（data augmentation）按任务约定变换或组合已有样本，以约束网络在相应输入上的预测。设随机变换为$T_{\xi}$，对应目标变换为$U_{\xi}$，增强后的目标为
\begin{equation}
  \hat R_{\mathrm{aug}}(\bm\theta)
  =\frac1n\sum_i\mathbb E_{\xi}
    [\ell(f_{\bm\theta}(T_{\xi}(\bm x_i)),U_{\xi}(y_i))]\text{。}
  \label{eq:generalization-augmentation-risk}
\end{equation}
分类中的保持标签变换令$U_{\xi}(y)=y$；分割中移动图像则需同步移动标签，表达输出随输入共同变换的等变关系。增强改变训练分布，而非增加与原始样本独立的新观测。

图像中的平移、裁剪、翻转和颜色调整，时间序列中的时间移位、局部扰动或伸缩，都需要按语义检查。例如，左右翻转不一定保持方向标签，裁剪可能移除识别对象，颜色变化可能破坏诊断所需信号，时间反转可能改变动作类别。变换名称本身不能证明标签保持不变。

增强强度决定约束覆盖的邻域。太弱时未必带来新的训练条件，太强时可能把样本推到目标定义已不适用的区域。训练时可以随机生成增强，验证时则按照确定的使用条件评估；若采用多次增强预测平均，应作为额外推理过程报告。仅从训练集估计的尺度、字典或变换参数也应保持数据边界，避免利用测试信息。

\subsection{基于样本混合的数据增强}
\label{sec:generalization-sample-mixing}

仅围绕各样本施加局部变换，仍可能让样本之间的预测任意变化。混合增强把插值或区域组合也交给同一个模型拟合，为这些位置规定额外监督关系。\term{Mixup}用输入和监督目标的凸组合训练模型。对输入$\bm x_i,\bm x_j$及独热类别向量$\bm y_i,\bm y_j$，令$\lambda\sim\operatorname{Beta}(a,a)$，$a>0$，构造
\begin{equation}
  \tilde{\bm x}=\lambda\bm x_i+(1-\lambda)\bm x_j,\qquad
  \tilde{\bm y}=\lambda\bm y_i+(1-\lambda)\bm y_j\text{。}
  \label{eq:generalization-mixup}
\end{equation}
交叉熵对目标向量线性，因此混合目标的损失等于$\lambda\ell(f(\tilde{\bm x}),\bm y_i)+(1-\lambda)\ell(f(\tilde{\bm x}),\bm y_j)$。这鼓励网络在样本之间变化得较为平缓，并将混合比例写进监督；它不声称混合图像是现实中的随机观测\scite{zhang2018mixup}

\term{CutMix}从一个样本保留部分区域，用另一个样本的对应区域替换。设二值空间掩码$\bm M$在保留区域为1，则$\tilde{\bm x}=\bm M\odot\bm x_i+(\bm1-\bm M)\odot\bm x_j$，目标仍用凸组合，但$\lambda$取实际保留面积占比。被裁到边界外的区域不应计入替换面积\scite{yun2019cutmix}

图~\ref{fig:generalization-mixing}用$2\times2$构造输入比较混合方式。Mixup使每个位置叠加，CutMix保留各区域的原始局部值。面积比例是分类监督的近似构造，不保证可见语义证据恰好按面积分配；若目标是逐像素分割，标签应按同一掩码组合，而非统一平均整图标签。
\begin{figure*}[htbp]
  \centering
  \begin{tikzpicture}[x=1mm,y=1mm,font=\figurefont\footnotesize]
  \foreach \offset/\title/\value in {0/样本A/1,42/样本B/0,84/Mixup/0.75,126/CutMix/1}{
    \begin{scope}[xshift=\offset mm]
      \node[nn heading] at (13,32) {\title};
      \foreach \i in {0,1}{\foreach \j in {0,1}{
        \def\cellvalue{\value}
        \ifnum\offset=126\relax
          \ifnum\i=1\relax\ifnum\j=1\relax\def\cellvalue{0}\fi\fi
        \fi
        \pgfmathsetmacro{\tone}{18+27*\cellvalue}
        \fill[NNInput!\tone] ({13*\j},{13-13*\i}) rectangle ({13+13*\j},{26-13*\i});
        \node[text=BookInk] at ({6.5+13*\j},{19.5-13*\i}) {$\cellvalue$};
      }}
      \draw[NNLine,line width=1.1pt] (0,0) rectangle (26,26);
      \draw[NNLine,line width=.5pt] (13,0)--(13,26) (0,13)--(26,13);
      \ifnum\offset=0\relax\node[nn annotation] at (13,-8) {$(1,0)^{\top}$};\fi
      \ifnum\offset=42\relax\node[nn annotation] at (13,-8) {$(0,1)^{\top}$};\fi
      \ifnum\offset>42\relax\node[nn annotation] at (13,-8) {$(0.75,0.25)^{\top}$};\fi
    \end{scope}
  }
\end{tikzpicture}

  \caption{构造的两类$2\times2$输入及其混合。样本A各位置为1，目标为$(1,0)^{\top}$；样本B各位置为0，目标为$(0,1)^{\top}$。Mixup取$\lambda=0.75$，各位置变为0.75；CutMix替换右下位置，保留面积也为0.75。二者目标均为$(0.75,0.25)^{\top}$，但空间组织不同。这是机制示例，不代表真实图像类别。}
  \label{fig:generalization-mixing}
\end{figure*}

混合强度、样本配对和目标构造应一起检查。若两个样本之间的插值没有任务意义，强行要求线性标签变化可能引入偏差；类别线索集中在很小区域时，面积权重也可能失真。方法的价值来自适当的函数约束，应通过保留原始语义的验证样本判断，而不是把所有混合都视为有效增加数据。

\section{模型侧的泛化控制}
\label{sec:generalization-constraints}

模型侧的控制规定允许哪些计算结构，以及用什么损失评价参数。参数共享与绑定直接缩小可表示的函数集合；范数惩罚保留参数范围，却在目标中增加尺度或稀疏偏好。两者分别对应第\ref{sec:generalization-capacity}节的归纳偏置与第\ref{sec:network-norm-bounds}节的复杂度分析。

这里按主要作用位置区分模型侧与训练侧：固定结构及显式损失项放在本节，更新收缩、随机训练扰动和轨迹选择放在第\ref{sec:generalization-training-process}节。方法在数学上可以联系，例如普通SGD中的L2梯度等价于特定收缩；分类仍需保留“改目标”与“改更新”的区别。

\subsection{参数共享与参数绑定}
\label{sec:generalization-parameter-sharing}

第\ref{sec:generalization-capacity}节的结构偏置可以直接体现为参数相等关系。共享减少需要分别估计的自由度，让多个位置共同提供梯度，但也把“这些位置适用同一规则”写进函数集合。同一种局部关系出现在多个位置时，分别学习所有参数会浪费样本。\term{参数共享}（parameter sharing）让不同计算位置使用同一参数对象。卷积核跨位置复用，循环权重跨时间复用；每次使用都贡献梯度，更新时累积到同一对象。\term{参数绑定}（weight tying）则强调多个模块的参数被限制为相等或具有指定对应关系，硬绑定常通过共享实现。

例如，词表嵌入矩阵$\bm E\in\mathbb R^{v\times d}$的第$j$行表示词元$j$，输出隐藏状态为$\bm h\in\mathbb R^d$时，可以绑定输出矩阵，令分数$\bm z=\bm E\bm h+\bm b$。这里要求输出隐藏宽度与嵌入宽度相容，输入和输出词表也有一致对应；不相容时需投影或另外定义绑定范围。绑定让同一词元的输入表示与输出判别共享信息，也减少独立参数，但同时施加两种职责可由同一表示承担的假设。

若不希望强制相等，可加入$\lambda\lVert\bm\theta_A-\bm\theta_B\rVert_2^2/2$，鼓励两组参数接近。有限$\lambda$下它们仍是两个参数对象，这是一种软约束。共享改变允许的函数结构，软约束改变目标中的偏好；即使参数量减少，也不意味着偏差必然降低。任务在不同位置或模块上需要不同规则时，过强绑定可能妨碍拟合。

\subsection{参数范数惩罚}
\label{sec:generalization-norm-penalties}

若允许的函数太丰富，目标只看训练拟合会留下很多参数选择。范数惩罚用额外代价偏好某些解，对应第\ref{sec:network-norm-bounds}节中的尺度复杂度；具体风险判断仍需连同输入和分类间隔。\term{范数惩罚}（norm penalty）将学习目标改为
\begin{equation}
  \min_{\bm\theta}\ \hat R_S(\bm\theta)+\lambda\Omega(\bm\theta),
  \qquad\lambda\geq0\text{。}
  \label{eq:generalization-regularized-objective}
\end{equation}
常见选择为$\Omega=\lVert\bm\theta\rVert_1$或$\Omega=\lVert\bm\theta\rVert_2^2/2$。严格说，后者是平方范数惩罚；实践中常简称L2正则化。它惩罚参数幅度，L1在0附近的非光滑形状则可以促成精确的零参数。

\begin{example}[L1与L2的收缩不同]
  \label{ex:generalization-shrinkage}
  对标量教学目标$\tfrac12(w-a)^2$，加入$\lambda w^2/2$后，最优解为$w=a/(1+\lambda)$；加入$\lambda|w|$后，最优解为$w=\operatorname{sign}(a)\max(|a|-\lambda,0)$。例如$a=0.4,\lambda=0.5$时，前者为$0.4/1.5$，后者为0。这个差异来自惩罚形状，不能据此认为任意普通梯度实现都会精确得到稀疏神经网络。
\end{example}

如果要一起选择或去掉一组参数，可以采用$\Omega(\bm\theta)=\sum_g\lVert\bm\theta_g\rVert_2$，其中组$g$可能对应某个特征、神经元或通道，必要时再按组大小加权。组范数对整组是否为0产生偏好，与逐参数L1的稀疏形态不同。重叠分组和实际模块结构需明确指定，不能只写“采用分组正则化”。

L1与组范数的零点不可微，近端更新常能更直接处理其非光滑部分；L2则有梯度$\lambda\bm\theta$。参数尺度还受重参数化影响，例如ReLU网络中某层放大、下一层相应缩小可能保留函数但改变参数惩罚。因此，参数范数不是完全独立于表示方式的函数复杂度。选择约束对象、系数及是否包含偏置和归一化参数，都应结合目标检查。

\section{训练侧的泛化控制}
\label{sec:generalization-training-process}

模型族和显式损失确定后，求解过程仍会影响最终预测。训练侧的方法管理何时停止、使用轨迹中的哪个参数、怎样更新，以及训练时施加什么随机扰动。它们联系第\ref{sec:algorithmic-stability}节的训练敏感性与第\ref{sec:implicit-bias}节的选解分析，评价重点是验证风险，而非继续压低训练误差。

下面先讨论停止、平均与优化器选择，再讨论权重衰减和随机扰动。权重衰减在更新中收缩参数，因此归入训练侧；Dropout及噪声按训练时的操作位置归类，也可以从随机计算的期望损失理解。后一联系说明分类并非互斥的数学定理，但有助于实现时明确该修改发生在哪里。

\subsection{提前终止}
\label{sec:generalization-early-stopping}

第\ref{sec:sgd-stability}节显示，在相应条件下训练时间和步长总量会进入稳定性分析。提前终止把“沿当前数据训练多久”作为选择变量，用验证表现限制继续拟合，而不是修改每一步的求导公式。持续训练可能继续降低训练误差，但验证表现已经不再改善。\term{提前终止}（early stopping）根据验证数据选择训练时刻，将训练轨迹上的参数位置当作候选模型。它需要规定评估指标、检查间隔、最小改善量和容忍窗口，并保存出现最佳验证表现时的参数。

例如，每训练一个轮次计算一次验证损失，若连续$p$次检查都没有超过阈值的改善，就停止并恢复已保存的最优参数。恢复参数与使用停止时刻的参数不同；继续训练还需相应恢复优化器和调度状态。随机验证增强、Dropout未关闭或BatchNorm模式错误，都可能让停止判断混入无关噪声。

图~\ref{fig:generalization-early-stopping}展示可能出现的曲线形状。真实验证曲线不一定具有单一谷底，短期变差也不证明后面不会改善。容忍窗口是等待与计算预算之间的选择，需要结合噪声、调度和目标，而非固定为普遍常数。
\begin{figure}[htbp]
  \centering
  \begin{tikzpicture}
  \begin{axis}[width=106mm,height=68mm,axis lines=left,
    xmin=0,xmax=80,ymin=0,ymax=.8,
    xlabel={训练轮次 $t$},ylabel={损失},
    xtick={0,20,40,60,80},ytick={0,.2,.4,.6,.8},
    tick label style={font=\figurefont\footnotesize},
    label style={font=\figurefont\footnotesize},
    legend style={draw=none,font=\figurefont\footnotesize,
      at={(.98,.98)},anchor=north east},domain=0:80,samples=161]
    \addplot[NNInput,line width=1.2pt] {.05+.7*exp(-x/15)};
    \addlegendentry{训练损失}
    \addplot[NNHidden,line width=1.2pt,dash pattern=on 3pt off 1.5pt] {.16+.5*exp(-x/12)+.002*x};
    \addlegendentry{验证损失}
    \draw[NNLine,dashed,line width=.85pt] (axis cs:36.43865,0)--(axis cs:36.43865,.8);
    \addplot[NNOutput,only marks,mark=*,mark size=2pt]
      coordinates {(36.43865,.2568773)};
  \end{axis}
\end{tikzpicture}

  \caption{提前终止的示意曲线。训练损失取$0.05+0.7\exp(-t/15)$，验证损失取$0.16+0.5\exp(-t/12)+0.002t$，$t$为训练轮次。虚线标出该构造验证函数的极小点附近，说明继续拟合与验证表现可能分离；这些函数仅用于讲解，不是实验记录。}
  \label{fig:generalization-early-stopping}
\end{figure}

训练时长限制了算法受训练数据影响的累计程度。在特定光滑、Lipschitz等条件下，SGD稳定性界可与步长总量及训练时间联系，第\ref{sec:sgd-stability}节给出了分析；这支持考察训练过程，而不能推出所有深度网络都越早停止越好\scite{hardt2016sgd}

验证集被用于选择停止时刻，反复尝试很多策略还可能对验证集过拟合。因此，最终测试应在策略确定后进行；时间序列和持续感知的验证还需模拟实际状态保留和未来信息边界。

\subsection{参数平均}
\label{sec:generalization-weight-averaging}

第\ref{sec:training-inexact-gradients}节中的更新噪声可能使后期参数在附近波动。平均这些位置尝试减弱某个瞬时状态的影响，以一次模型推理综合同一轨迹的信息，但平均后是否仍处在合适损失区域需要检查。\term{指数移动平均}（exponential moving average, EMA）维护$\bar{\bm\theta}^{(t)}=\rho\bar{\bm\theta}^{(t-1)}+(1-\rho)\bm\theta^{(t)}$，$\rho\in[0,1)$。取初值为训练初始参数或当前参数可以避免零初值的缩小效应，$\rho$越大，对近期变化的响应越慢。推理时使用平均参数形成一个模型。

\term{随机权重平均}（stochastic weight averaging, SWA）选取训练后段的多个参数状态作算术平均。对已收集的$k$个状态，加入新状态$\bm\theta$时更新$\bar{\bm\theta}_{k+1}=k\bar{\bm\theta}_k/(k+1)+\bm\theta/(k+1)$。原始方案利用较大或周期性学习率采样训练后段区域，平均后若使用BatchNorm，还需按平均模型重新估计运行统计量\scite{izmailov2018swa}

参数平均与预测平均不同。非线性网络通常不满足$f_{(\bm\theta_A+\bm\theta_B)/2}=(f_{\bm\theta_A}+f_{\bm\theta_B})/2$。对同一训练轨迹附近、坐标含义一致的参数，平均可能仍落在具有良好损失的区域；对独立模型，则可能因隐藏单元的置换或尺度不一致而失败。

\begin{example}[等价网络的参数平均可以失效]
  \label{ex:generalization-parameter-average}
  取无偏置的两单元ReLU网络$f(x)=\operatorname{ReLU}(x)+\operatorname{ReLU}(-x)=|x|$。网络A的输入权重为$(1,-1)$，输出权重为$(1,1)$；网络B只交换两个隐藏单元，输入权重为$(-1,1)$，输出权重仍为$(1,1)$。两个网络的函数完全相同，但逐坐标平均后的输入权重为$(0,0)$，输出恒为0。预测平均仍给出$|x|$。
\end{example}

平均适用性应通过平均参数的实际验证表现判断。EMA和SWA减少了保存全部预测模型的需要，但不保证不增加偏差，也不等同于对独立训练随机性的完整集成。

\subsection{优化算法的隐式偏置}
\label{sec:generalization-implicit-bias}

同一目标可能有多个低损失解，更新方向决定训练倾向哪一个。\term{隐式偏置}（implicit bias，前文亦称隐式偏差）指学习过程在这些解中形成的选择偏好，与小批量梯度是否无偏是不同概念。第\ref{sec:implicit-bias}节的线性与矩阵例子说明，这种偏好可以在没有额外惩罚项时出现。

可分线性二分类提供一个明确结果：对$y_i\in\{-1,1\}$、无显式惩罚的逻辑损失，在相应步长与可分条件下，梯度下降使$\lVert\bm w\rVert_2$不断增长，方向趋向欧氏最大间隔分隔方向。损失没有有限范数的零风险极小点，更新过程仍选择了特定方向；其条件与推导见第\ref{sec:implicit-bias}节\scite{soudry2018implicit}

\begin{example}[改变更新几何会改变选出的解]
  \label{ex:generalization-update-geometry}
  用一个训练样本拟合$w_1+2w_2=1$，目标为$\tfrac12(w_1+2w_2-1)^2$。从零初始化出发，足够小的固定步长梯度下降始终沿$(1,2)^{\top}$移动，收敛到$(1/5,2/5)^{\top}$，这是欧氏最小范数解。若改为固定预条件$\bm P=\operatorname{diag}(4,1)$，更新始终沿$(4,2)^{\top}$移动，则收敛到$(1/2,1/4)^{\top}$。两个解的训练误差都为0，后者最小化的是$\bm w^{\top}\bm P^{-1}\bm w$。未来数据若需要不同坐标承担不同作用，二者的风险也会不同；仅凭这一个训练约束不能确定哪个偏好正确。
\end{example}

深度网络中的偏好还随参数化、初始化、学习率、批量与训练时长改变，因此实践不是直接宣布某个优化器“泛化最好”。Wilson等给出了自适应方法与GD／SGD选出不同解的构造，并在所研究的网络中观察到自适应方法训练更好而测试更差。这提供了比较非自适应与自适应更新的理由，结论限于相应设置\scite{wilson2017adaptive}

Choi等进一步发现，优化器排名对超参数搜索范围和调节预算敏感。因而把调节充分的SGD与默认Adam比较，或只比较相同学习率数值，都不足以判断算法的泛化优势。AdamW还含解耦衰减，比较时需要把更新规则和显式收缩一起说明\scite{choi2020optimizers}

把隐式偏置用于训练选择，可以采用以下可执行的比较过程：
\begin{enumerate}
  \item 固定网络、数据划分、增强与损失，以带动量SGD和AdamW作为两种候选。已有训练方案能稳定拟合时，先保留其为基准；更换优化器是检验不同更新几何的一种方式。
  \item 分别调节学习率、调度、动量或矩系数，以及适用的衰减系数。给两类候选相当的搜索与训练预算，记录达到目标训练损失的成本和验证表现；不要强制共享同一组数值。
  \item 联合阅读训练与验证曲线。训练误差仍高时，先排查第\ref{chap:neural-network-training}章的优化困难；训练相近而验证明显不同，才有较明确的选解差异。固定时间预算的表现也应单独报告，因为慢方法可能尚未完成拟合。
  \item 对有希望的配置使用多个随机种子复核，再用验证结果选择。停止时刻、批量和参数平均也会改变轨迹，可以分别作控制实验；最终测试集留到方案确定之后。
\end{enumerate}

如果某一配置持续取得更好的验证风险，就在该任务上采用它；目前这些理论与实验没有给出对所有网络都更好的一种优化器。训练中途从AdamW切到SGD也是新的训练方案，需要重新设置学习率与状态处理并验证，不能默认保留旧矩状态或保证切换后泛化更好。

选解指标本身也要有边界。“平坦极小点泛化更好”需要指定参数坐标与扰动尺度，保持函数不变的重缩放可以改变曲率。使用平坦性评价时，应说明怎样处理这种歧义，不能只凭损失曲面看起来宽就判定未来风险更低\scite{dinh2017sharp}

\subsection{权重衰减}
\label{sec:generalization-weight-decay}

第\ref{sec:training-optimizers}节把方向、尺度和历史统计分开，参数收缩也可以独立加入更新。权重衰减通过这种算法选择控制尺度；是否等价于某个惩罚目标，要检查预条件和历史缓冲。范数惩罚改变目标，\term{权重衰减}（weight decay）直接在更新中收缩参数。对无动量的普通SGD，L2惩罚更新为
\begin{equation}
  \bm\theta^{(t)}=\bm\theta^{(t-1)}-
    \eta_t(\bm g^{(t)}+\lambda\bm\theta^{(t-1)})
  =(1-\eta_t\lambda)\bm\theta^{(t-1)}-\eta_t\bm g^{(t)}\text{。}
  \label{eq:generalization-sgd-weight-decay}
\end{equation}
因而在相同参数范围及系数约定下，L2惩罚与学习率乘以$\lambda$的收缩等价。若所谓“衰减系数”直接表示每步乘法系数，就还需对应学习率；加入动量也要检查惩罚是否进入动量缓冲，不能不看实现就延用等价结论。

对预条件矩阵$\bm P^{(t)}$，惩罚项进入梯度后产生$-\eta_t\lambda\bm P^{(t)}\bm\theta$，而解耦收缩产生$-\eta_t\lambda\bm\theta$。二者通常不同，Adam中惩罚还可能进入一阶、二阶矩统计。\term{AdamW}将衰减从自适应梯度更新中分离，其基本更新为
\begin{equation}
  \bm\theta^{(t)}=(1-\eta_t\lambda)\bm\theta^{(t-1)}-
    \eta_t\frac{\hat{\bm m}^{(t)}}{\sqrt{\hat{\bm v}^{(t)}}+\epsilon}\text{，}
  \label{eq:generalization-adamw}
\end{equation}
这里矩统计来自数据损失梯度，具体计算见式\eqref{eq:training-adam}\scite{loshchilov2019adamw}

衰减通过多步乘积发生，学习率调度、更新次数和是否排除部分参数都会改变总收缩量。普通参数坐标下的解耦权重衰减，也不等同于所有预条件坐标中的同一L2目标。比较正则化强度时，需要固定更新规则和系数含义。

\subsection{Dropout}
\label{sec:generalization-dropout}
\label{sec:generalization-stochastic-regularization}

当低训练误差依赖一组必须共同出现的中间特征时，可以改变训练时可用的计算路径，要求共享参数在随机子网络中继续预测。Dropout因此约束受扰动的模型行为，其直接目标不是让单次无扰动损失下降更快。\term{Dropout}在训练时随机屏蔽部分激活，使网络在不同的可用特征子集上完成同一任务。设隐藏表示为$\bm h$，各元素独立保留概率为$q\in(0,1]$，$m_j\sim\operatorname{Bernoulli}(q)$。常见的反向缩放形式为
\begin{equation}
  \tilde h_j=\frac{m_j}{q}h_j,\qquad
  \mathbb E[\tilde h_j\mid\bm h]=h_j,\qquad
  \operatorname{Var}(\tilde h_j\mid\bm h)=\frac{1-q}{q}h_j^2\text{。}
  \label{eq:generalization-dropout}
\end{equation}
训练时除以$q$，常规推理时关闭屏蔽并直接使用$\bm h$。另一约定在训练时不缩放、推理时乘$q$，两者必须与实现一致。这里$q$为保留率，有些接口使用的是失活率$1-q$\scite{srivastava2014dropout}

训练目标包含对掩码的期望$\mathbb E_{\bm m}\ell(f_{\bm\theta,\bm m}(\bm x),y)$。随机子网络共享参数，它们并不是独立训练的多个模型。保持隐藏表示的期望，也不意味着完整非线性网络的输出满足$\mathbb E_{\bm m}f_{\bm\theta,\bm m}=f_{\bm\theta}$；因此，单次无掩码推理是某种模型平均的近似，不能当成精确集成公式。

图~\ref{fig:generalization-dropout}展示同一隐藏层的一个掩码样本。屏蔽后其输入参数仍保留，后续批量可能重新启用，并非永久剪枝。较低$q$增加扰动和方差，也可能损害拟合。BatchNorm与Dropout组合时，还需关注随机激活对运行统计的影响；序列模型中跨时间独立采样或共享掩码，则具有不同扰动结构。
\begin{figure}[htbp]
  \centering
  \begin{tikzpicture}[x=1mm,y=1mm,font=\figurefont\footnotesize]
  \foreach \i/\y in {1/12,2/-12}{\node[nn input] (in\i) at (8,\y) {$x_\i$};}
  \node[nn hidden] (h1) at (48,23) {$h_1/q$};
  \node[nn neuron,draw=NNLine,fill=NNCanvas,text=BookMuted,dashed] (h2) at (48,0) {$0$};
  \node[nn hidden] (h3) at (48,-23) {$h_3/q$};
  \node[nn output] (out) at (88,0) {$y$};
  \foreach \i in {1,2}{
    \foreach \j in {1,3}{\draw[nn link] (in\i)--(h\j);}
    \draw[nn link,dashed] (in\i)--(h2);
  }
  \foreach \j in {1,3}{\draw[nn flow,draw=NNOutput] (h\j)--(out);}
  \draw[nn link,dashed] (h2)--(out);
  \draw[NNLine,line width=1.1pt] (45,-3)--(51,3) (45,3)--(51,-3);
  \node[nn heading] at (8,35) {输入};
  \node[nn heading] at (48,35) {随机失活};
  \node[nn heading] at (88,35) {输出};
\end{tikzpicture}

  \caption{一个三单元隐藏层的构造性Dropout示意。中间单元被屏蔽，灰色虚线为本次不参与前向计算的路径，其余实线仍传递信息。此图只展示一个掩码样本；实际保留概率不由图中剩余单元比例决定，训练时按式\eqref{eq:generalization-dropout}缩放。}
  \label{fig:generalization-dropout}
\end{figure}

\subsection{结构化随机失活}
\label{sec:generalization-structured-dropout}

随机扰动应覆盖模型真正依赖的信息单位；若重复线索能从邻近位置轻易补回，逐元素屏蔽施加的约束就较弱。结构化失活按连接、通道、块或层改变有效路径，与第\ref{sec:training-connections}节的结构组织相对应。邻近位置的特征可能高度相关，独立屏蔽某个像素位置的激活后，其信息仍能由邻近位置提供。结构化随机失活根据网络结构选择更大或不同的屏蔽对象，而非只提高逐元素失活概率。

\term{DropConnect}对权重连接采样掩码，令$\tilde{\bm W}=\bm M\odot\bm W$，随后再计算激活；Dropout则直接作用于激活。基本训练与推理的缩放、分布近似应按各自约定处理，两种失活的随机函数不相同\scite{wan2013dropconnect}

\term{空间Dropout}（spatial dropout）在一个通道的所有空间位置使用同一掩码，削弱网络对某个整通道特征的依赖。\term{DropBlock}随机屏蔽特征图中连续的空间块，干扰局部邻域可提供的重复信息。块大小、采样位置和边界会影响实际失活率，不能仅把块中心采样概率当作元素失活概率\scite{ghiasi2018dropblock}

\term{随机深度}（stochastic depth）对残差分支随机启用或绕过，让训练时经过的有效深度变化。一个期望保持的实现写为$\bm h^{(l+1)}=\bm h^{(l)}+(m_l/q_l)F_l(\bm h^{(l)})$，$m_l\sim\operatorname{Bernoulli}(q_l)$，推理时使用全部分支。原始方法也可将缩放放在推理阶段；是否真正跳过被屏蔽分支的计算影响训练时间\scite{huang2016stochasticdepth}

连接、通道、空间块和整层对应不同的信息粒度。屏蔽过大可能毁掉任务所需信息，特别是对象很小、通道很窄或缺少恒等路径时。失活应与结构匹配，并报告保留概率、作用位置及采样共享范围。

\subsection{噪声注入}
\label{sec:generalization-noise}

第\ref{sec:parameter-distributions}节把参数分布及其期望损失纳入分析；噪声训练也可以看成在规定扰动分布下优化平均表现。其作用由注入位置和协方差决定，不能只凭“加入随机性”推断泛化改善。随机扰动还可以连续改变表示或参数，而非只选择保留与屏蔽。\term{噪声注入}（noise injection）在训练中使用$\tilde{\bm h}=\bm h+\bm\epsilon$或$\tilde{\bm\theta}=\bm\theta+\bm\epsilon$，其中噪声的均值、协方差和采样范围需要明确。输入噪声属于数据增强的一种，中间表示噪声改变后续计算，参数噪声则同时影响所有使用该参数的位置。

对足够光滑的目标函数和小幅零均值噪声，二阶展开给出局部近似
\begin{equation}
  \mathbb E_{\bm\epsilon}[\mathcal L(\bm\theta+\bm\epsilon)]
  \approx\mathcal L(\bm\theta)+\tfrac12
      \operatorname{tr}(\bm\Sigma\nabla^2\mathcal L(\bm\theta))\text{，}
  \label{eq:generalization-noise-expansion}
\end{equation}
其中$\bm\Sigma=\mathbb E[\bm\epsilon\bm\epsilon^{\top}]$。近似的余项需要噪声幅度与光滑条件控制；二阶项未必非负，所以不能把任意非凸目标上的噪声都解释为固定的正范数惩罚。

一个可精确核查的特殊情况是线性回归。对预测$\bm w^{\top}\bm x$、平方损失及输入噪声$\mathbb E\bm\epsilon=\bm0$、$\mathbb E\bm\epsilon\bm\epsilon^{\top}=\sigma^2\bm I$，有
\begin{equation}
  \mathbb E\tfrac12[y-\bm w^{\top}(\bm x+\bm\epsilon)]^2
  =\tfrac12(y-\bm w^{\top}\bm x)^2+\tfrac{\sigma^2}{2}\lVert\bm w\rVert_2^2\text{。}
  \label{eq:generalization-linear-noise}
\end{equation}
零均值消去交叉项，协方差决定额外二次项。这说明噪声与正则化的联系取决于损失、模型和注入位置；非线性网络中的输入噪声通常不能直接化为所有参数的L2惩罚。

Dropout也可理解为一种乘性噪声：$\tilde h_j=h_j+[m_j/q-1]h_j$，条件均值为0、方差与$h_j^2$成比例。它与加性高斯噪声的分布和尺度不同。噪声是否符合任务中合理的变化，决定了对局部邻域的约束是否有意义；它不能修复系统性错误标签或数据泄漏。

\section{模型集成}
\label{sec:generalization-ensembles}

若不同训练过程保留了不同预测误差，可以通过组合预测减少其中部分波动。这个机制与第\ref{sec:generalization-solution-selection}节的解差异相连；收益取决于成员质量与误差相关性，而不是单纯增加模型数。

\subsection{Bagging}
\label{sec:generalization-bagging}

Bagging围绕训练样本变化生成多个解，用预测组合减弱单次采样的影响。它针对的是学习过程的抽样敏感性，与第\ref{sec:algorithmic-stability}节关注的问题相邻，不过重采样平均本身不是该稳定性定理的证明。\term{Bagging}（bootstrap aggregating，自助聚合）从$n$个训练样本中有放回抽取$n$次构成一个重采样集，对多个这样的数据集分别训练模型，回归常用预测平均，分类可用概率平均或多数投票\scite{breiman1996bagging}

对第$m$个分类模型的概率向量$\bm p_m(\bm x)$，$M$个成员的平均为$\bar{\bm p}(\bm x)=M^{-1}\sum_m\bm p_m(\bm x)$，预测类别取其最大分量。概率平均与类别投票不同，概率平均也通常不等于先平均分数再做Softmax。组合方式属于预测规则，应在评价前固定。

重采样集含重复记录，一个原始样本被遗漏的概率为$(1-1/n)^n$，不同原始样本的期望覆盖数为$n[1-(1-1/n)^n]$，大$n$时约为$0.632n$。未被某成员使用的样本可用于该成员的袋外评估，但神经网络的超参数选择仍需保持评估边界。相关数据需要按主体或时间块重采样；逐行抽样不能模拟独立主体变化。

集成收益与成员误差相关性有关。固定某个输入，假设各成员误差的方差均为$\sigma^2$，不同成员的协方差均为$\rho\sigma^2$，则平均误差的方差为
\begin{equation}
  \operatorname{Var}\left(\frac1M\sum_{m=1}^{M}e_m\right)
  =\sigma^2\left[\rho+\frac{1-\rho}{M}\right]\text{。}
  \label{eq:generalization-ensemble-variance}
\end{equation}
这是由$M$个方差项和$M(M-1)$个协方差项求和得到的。独立误差时方差按$1/M$下降，完全同向误差时不下降；共同偏差也不会因平均自动消失。这个等方差、等相关的教学模型用于解释关系，真实成员未必满足这些条件。

\subsection{基于训练随机性的集成}
\label{sec:generalization-randomness-ensemble}

第\ref{sec:generalization-solution-selection}节中，初始化与数据次序即使在同一训练集上也能改变所得解。利用这类差异构造成员，可以不重采样原始数据，但仍需检查误差是否足够不同。重采样数据只是获得不同成员的一种方式。也可以让所有网络使用同一训练集，通过不同初始化、数据顺序、增强和随机失活形成不同预测函数。这与Bagging的数据组织不同：训练随机性集成不要求对原始样本进行有放回重采样，也不自然提供相同的袋外集合。

成员的质量与差异需要一起考虑。加入一个在关键输入上系统性错误的模型，未必改善平均；只改变随机种子但各模型预测几乎一致，也可能收效有限。应在同一验证集上考察单模型表现与预测分歧，并固定组合权重；不能通过查看最终测试结果反复筛选成员。

集成通常增加训练和推理成本。$M$个相同规模网络若顺序运行，模型计算量约为单个的$M$倍；并行运行可以降低部分等待时间，却需要更多设备或存储。参数平均形成一个模型，预测集成保留多个模型，这个区别决定部署成本。选择成员数与组合方式时，任务质量、延迟和资源预算都应进入判断。

\section{本章小结}
\label{sec:generalization-summary}

低训练误差不能单独说明网络学得了可迁移规律。未来风险受到数据覆盖与质量、模型归纳偏置和训练解选择共同影响。同分布评估、场景变化和持续状态的边界不同，数据划分必须对应预期使用方式。

增强把任务允许的变化写进经验目标，参数共享与范数惩罚引入结构和尺度偏好，权重衰减通过更新收缩参数，随机失活与噪声约束受扰动的计算，停止时刻和参数平均改变轨迹上的解选择，集成则利用多个预测的差异。它们可以组合，但不具有无条件的收益：错误增强会改变语义，过强约束会损害拟合，不对齐的参数平均会破坏函数，高度相关的模型也难以从集成获益。将这些机制与复杂度、稳定性和隐式偏置联系起来，既有助于解释观察，也要求保留各自的条件；最终选择仍应由符合使用场景的验证和独立测试支持。
